%%%%%%%%%%%%%%%%
%%% num 11 %%%
%%%%%%%%%%%%%%%%

\Chapter{11}

\Verse{1}
{
	\hE{וַיְהִי הָעָם כְּמִתְאֹנְנִים רַע בְּאזְנֵי יְהֹוָה וַיִּשְׁמַע יְהֹוָה וַיִּחַר אַפּוֹ וַתִּבְעַר בָּם אֵשׁ יְהֹוָה וַתֹּאכַל בִּקְצֵה הַמַּחֲנֶה׃}
}
{
	\eE{
		And the people were like grumblers, bad in YHWH’s ears,
		and YHWH heard, and His anger flared, and a fire of YHWH
		burned among them and consumed at the edge of the camp.
	}
}
{
%	\footnote{
%		the people were like grumblers. This is just the first of many accounts
%		of their complaining in Numbers, and, interestingly, it is the only one
%		in which the reason for their dissatisfaction is not told. The story’s
%		function, therefore, is not to point to any of the obstacles in the
%		wilderness but rather to indicate the people’s state of mind: negative,
%		volatile, unconfident.
%	}
}

\Verse{2}
{
	\hE{וַיִּצְעַק הָעָם אֶל מֹשֶׁה וַיִּתְפַּלֵּל מֹשֶׁה אֶל יְהֹוָה וַתִּשְׁקַע הָאֵשׁ׃}
}
{
	\eE{
		And the people cried out to Moses, and Moses prayed to
		YHWH, and the fire subsided.
	}
}
{
%	\footnote{
%		the people cried out to Moses. The people appeal to Moses rather than to
%		God, and Moses acts as intercessor. They still turn to the human leader
%		rather than directly to the mysterious, invisible divine source.
%	}
}

\Verse{3}
{
	\hE{וַיִּקְרָא שֵׁם הַמָּקוֹם הַהוּא תַּבְעֵרָה כִּי בָעֲרָה בָם אֵשׁ יְהֹוָה׃}
}
{
	\eE{
		And he called that place’s name Taberah because a fire of
		YHWH burned among them.
	}
}
{
}

\Verse{4}
{
	\hE{וְהָאסַפְסֻף אֲשֶׁר בְּקִרְבּוֹ הִתְאַוּוּ תַּאֲוָה וַיָּשֻׁבוּ וַיִּבְכּוּ גַּם בְּנֵי יִשְׂרָאֵל וַיֹּאמְרוּ מִי יַאֲכִלֵנוּ בָּשָׂר׃}
}
{
	\eE{
		And the gathered mass who were among them had a longing,
		and the children of Israel, as well, went back and cried,
		and they said, “Who will feed us meat?
	}
}
{
%	\footnote{
%		the gathered mass. See the comment on Num 12:2.
%	}
}

\Verse{5}
{
	\hE{זָכַרְנוּ אֶת הַדָּגָה אֲשֶׁר נֹאכַל בְּמִצְרַיִם חִנָּם אֵת הַקִּשֻּׁאִים וְאֵת הָאֲבַטִּחִים וְאֶת הֶחָצִיר וְאֶת הַבְּצָלִים וְאֶת הַשּׁוּמִים׃}
}
{
	\eE{
		We remember the fish that we would eat in Egypt for free:
		the cucumbers and the melons and the leek and the onions
		and the garlics.
	}
}
{
}

\Verse{6}
{
	\hE{וְעַתָּה נַפְשֵׁנוּ יְבֵשָׁה אֵין כֹּל בִּלְתִּי אֶל הַמָּן עֵינֵינוּ׃}
}
{
	\eE{
		And now, our soul is dried up. There isn’t anything—except
		the manna before our eyes.”
	}
}
{
}

\Verse{7}
{
	\hE{וְהַמָּן כִּזְרַע גַּד הוּא וְעֵינוֹ כְּעֵין הַבְּדֹלַח׃}
}
{
	\eE{
		And the manna: it was like a seed of coriander, and its
		appearance was like the appearance of bdellium.
	}
}
{
}

\Verse{8}
{
	\hE{שָׁטוּ הָעָם וְלָקְטוּ וְטָחֲנוּ בָרֵחַיִם אוֹ דָכוּ בַּמְּדֹכָה וּבִשְּׁלוּ בַּפָּרוּר וְעָשׂוּ אֹתוֹ עֻגוֹת וְהָיָה טַעְמוֹ כְּטַעַם לְשַׁד הַשָּׁמֶן׃}
}
{
	\eE{
		The people went around and collected and ground it in
		mills or pounded it in a mortar and cooked it in a pot and
		made it into cakes. And its taste was like the taste of
		something creamy made with oil.
	}
}
{
}

\Verse{9}
{
	\hE{וּבְרֶדֶת הַטַּל עַל הַמַּחֲנֶה לָיְלָה יֵרֵד הַמָּן עָלָיו׃}
}
{
	\eE{
		And when the dew descended on the camp at night, the manna
		would descend on it.
	}
}
{
}

\Verse{10}
{
	\hE{וַיִּשְׁמַע מֹשֶׁה אֶת הָעָם בֹּכֶה לְמִשְׁפְּחֹתָיו אִישׁ לְפֶתַח אהֳלוֹ וַיִּחַר אַף יְהֹוָה מְאֹד וּבְעֵינֵי מֹשֶׁה רָע׃}
}
{
	\eE{
		And Moses heard the people crying by their families, each
		at his tent entrance, and YHWH’s anger flared very much,
		and it was bad in Moses’ eyes.
	}
}
{
}

\Verse{11}
{
	\hE{וַיֹּאמֶר מֹשֶׁה אֶל יְהֹוָה לָמָה הֲרֵעֹתָ לְעַבְדֶּךָ וְלָמָּה לֹא מָצָתִי חֵן בְּעֵינֶיךָ לָשׂוּם אֶת מַשָּׂא כּל הָעָם הַזֶּה עָלָי׃}
}
{
	\eE{
		And Moses said to YHWH, “Why have you done bad to your
		servant, and why have I not found favor in your eyes, to
		set the burden of this entire people on me?
	}
}
{
%	\footnote{
%		Why have you done bad to your servant. Is this a way to talk to God?! In
%		what is the most extraordinary speech spoken by a human to God in the
%		Torah, Moses begins with these words and ends with: “And if this is how
%		you treat me, kill me!” We have seen a dramatic change in the way humans
%		speak to God: Adam and Cain speak like children caught in a naughty act.
%		Abraham dares to question a divine decision (the destruction of Sodom)
%		and says, “Far be it from you to do a thing like this.... Will the judge
%		of all the earth not do justice?” (Gen 18:25). Moses pleads for the
%		people after the golden calf event and says that if God will not forgive
%		them, then “wipe me out from your scroll that you’ve written” (Exod
%		32:32). And now Moses, pleading for himself, speaks to the creator in a
%		way that we would never imagine Adam or Abraham or Jacob or even Moses
%		himself speaking until now. It can only be understood in the light of
%		everything that has happened up to this moment: God’s speaking to Moses
%		in a personal way from the beginning (“I’m your father’s God”), Moses’
%		spending long hours in conversation with God, Moses’ performance of the
%		miracles. The same Moses whose first reaction to the deity was to “hide
%		his face, because he was afraid of looking at God” (Exod 3:6) now dares
%		to speak forcefully. It is the opposite of when an atheist uses some
%		audacious or outrageous language about gods. Moses is pictured as the
%		human who has come the closest to God, the one person who has seen God,
%		and it is precisely Moses who speaks the most audaciously of all. And
%		Moses in turn thus epitomizes humankind: it is precisely the wilderness
%		generation, the generation that is closest to God, who hear God’s voice
%		and see miracles daily, who are the most rebellious generation in the
%		Bible. This in turn may remind us of the Bible’s first story: when
%		humans live in the garden where God walks, they make the initial act of
%		rebellion. An ongoing point of the Torah is thus that humans and their
%		creator have an underlying basis of conflict, something in the very
%		nature of the divine-human condition. And the relations between God and
%		Abraham’s descendants are aimed at eventually resolving this. Thus Moses
%		is not seen as bad for speaking this way to God, and he is not punished.
%		It is rather part of his close relationship with God, and part of the
%		Torah’s central story.
%	}
}

\Verse{12}
{
	\hE{הֶאָנֹכִי הָרִיתִי אֵת כּל הָעָם הַזֶּה אִם אָנֹכִי יְלִדְתִּיהוּ כִּי תֹאמַר אֵלַי שָׂאֵהוּ בְחֵיקֶךָ כַּאֲשֶׁר יִשָּׂא הָאֹמֵן אֶת הַיֹּנֵק עַל הָאֲדָמָה אֲשֶׁר נִשְׁבַּעְתָּ לַאֲבֹתָיו׃}
}
{
	\eE{
		Did I conceive this entire people? Did I give birth to it,
		that you should say to me, ‘Carry it in your bosom,’ the
		way a nurse carries a suckling, to the land that you swore
		to its fathers?
	}
}
{
%	\footnote{
%		by myself, to carry. Moses complains that he has to bear the burden of
%		the people alone. This is in explicit contrast to the report about
%		Jethro’s advice to Moses in Exod 18:22, in which Jethro tells him to
%		pick judges to help him with the weight of his task, “and they’ll bear
%		it with you.” Moses takes Jethro’s advice there and appoints judges to
%		help him, so why does he say now that he is carrying the burden alone?
%		Notably, that earlier matter did not come from God, and the text never
%		says that God sanctioned it. It was specifically Moses’ own arrangement,
%		on the advice of his father-in-law. He still complains now to God that
%		God has left so much to him.
%	}
}

\Verse{13}
{
	\hE{מֵאַיִן לִי בָּשָׂר לָתֵת לְכל הָעָם הַזֶּה כִּי יִבְכּוּ עָלַי לֵאמֹר תְּנָה לָּנוּ בָשָׂר וְנֹאכֵלָה׃}
}
{
	\eE{
		From where do I have meat to give to this entire people,
		that they cry at me, saying, ‘Give us meat, and let’s
		eat’?
	}
}
{
}

\Verse{14}
{
	\hE{לֹא אוּכַל אָנֹכִי לְבַדִּי לָשֵׂאת אֶת כּל הָעָם הַזֶּה כִּי כָבֵד מִמֶּנִּי׃}
}
{
	\eE{
		I’m not able, I, by myself, to carry this entire people,
		because it’s too heavy for me.
	}
}
{
%	\footnote{
%		by myself, to carry. Moses complains that he has to bear the burden of
%		the people alone. This is in explicit contrast to the report about
%		Jethro’s advice to Moses in Exod 18:22, in which Jethro tells him to
%		pick judges to help him with the weight of his task, “and they’ll bear
%		it with you.” Moses takes Jethro’s advice there and appoints judges to
%		help him, so why does he say now that he is carrying the burden alone?
%		Notably, that earlier matter did not come from God, and the text never
%		says that God sanctioned it. It was specifically Moses’ own arrangement,
%		on the advice of his father-in-law. He still complains now to God that
%		God has left so much to him. Footnote 11:14. too heavy for me. The
%		recurring adjective “heavy,” which persisted through the exodus account
%		and culminated with the heavy cloud and the divine glory (kbôd),
%		now returns, with Moses using it to describe the impossible burden on
%		him, just as Jethro had told Moses earlier (Exod 18:18). Note the
%		parallels in Jethro’s words there and Moses’ words here, which indicate
%		that Moses’ father-in-law indeed had much influence on him. Moses’ use
%		of the term “heavy” now recalls the exodus, arguably to remind God of
%		Moses’ past service.
%	}
}

\Verse{15}
{
	\hE{וְאִם כָּכָה אַתְּ עֹשֶׂה לִּי הרְגֵנִי נָא הָרֹג אִם מָצָאתִי חֵן בְּעֵינֶיךָ וְאַל אֶרְאֶה בְּרָעָתִי׃}
}
{
	\eE{
		And if this is how you treat me, kill me, if I’ve found
		favor in your eyes, and let me not see my suffering.”
	}
}
{
}

\Verse{16}
{
	\hE{וַיֹּאמֶר יְהֹוָה אֶל מֹשֶׁה אֶסְפָה לִּי שִׁבְעִים אִישׁ מִזִּקְנֵי יִשְׂרָאֵל אֲשֶׁר יָדַעְתָּ כִּי הֵם זִקְנֵי הָעָם וְשֹׁטְרָיו וְלָקַחְתָּ אֹתָם אֶל אֹהֶל מוֹעֵד וְהִתְיַצְּבוּ שָׁם עִמָּךְ׃}
}
{
	\eE{
		And YHWH said to Moses, “Gather to me seventy men from
		Israel’s elders whom you’ve known because they’re the
		people’s elders and its officers, and you’ll take them to
		the Tent of Meeting, and they’ll stand up there with you.
	}
}
{
}

\Verse{17}
{
	\hE{וְיָרַדְתִּי וְדִבַּרְתִּי עִמְּךָ שָׁם וְאָצַלְתִּי מִן הָרוּחַ אֲשֶׁר עָלֶיךָ וְשַׂמְתִּי עֲלֵיהֶם וְנָשְׂאוּ אִתְּךָ בְּמַשָּׂא הָעָם וְלֹא תִשָּׂא אַתָּה לְבַדֶּךָ׃}
}
{
	\eE{
		And I’ll come down and speak with you there, and I’ll take
		some of the spirit that is on you, and I’ll set it on
		them, and they’ll carry the burden of the people with you,
		and you won’t bear it, you, by yourself.
	}
}
{
}

\Verse{18}
{
	\hE{וְאֶל הָעָם תֹּאמַר הִתְקַדְּשׁוּ לְמָחָר וַאֲכַלְתֶּם בָּשָׂר כִּי בְּכִיתֶם בְּאזְנֵי יְהֹוָה לֵאמֹר מִי יַאֲכִלֵנוּ בָּשָׂר כִּי טוֹב לָנוּ בְּמִצְרָיִם וְנָתַן יְהֹוָה לָכֶם בָּשָׂר וַאֲכַלְתֶּם׃}
}
{
	\eE{
		And you shall say to the people, ‘Consecrate yourselves
		for tomorrow, and you’ll eat meat, because you cried in
		YHWH’s ears saying: Who will feed us meat, because we had
		it good in Egypt. And YHWH will give you meat, and you’ll
		eat.
	}
}
{
}

\Verse{19}
{
	\hE{לֹא יוֹם אֶחָד תֹּאכְלוּן וְלֹא יוֹמָיִם וְלֹא חֲמִשָּׁה יָמִים וְלֹא עֲשָׂרָה יָמִים וְלֹא עֶשְׂרִים יוֹם׃}
}
{
	\eE{
		You’ll eat not for one day and not two days and not five
		days and not ten days and not ten days and not twenty
		days.
	}
}
{
}

\Verse{20}
{
	\hE{עַד חֹדֶשׁ יָמִים עַד אֲשֶׁר יֵצֵא מֵאַפְּכֶם וְהָיָה לָכֶם לְזָרָא יַעַן כִּי מְאַסְתֶּם אֶת יְהֹוָה אֲשֶׁר בְּקִרְבְּכֶם וַתִּבְכּוּ לְפָנָיו לֵאמֹר לָמָּה זֶּה יָצָאנוּ מִמִּצְרָיִם׃}
}
{
	\eE{
		Until a month of days! Until it will come out of your
		nose, and it will be a revulsion to you! Because you
		rejected YHWH, who was among you, and you cried in front
		of him, saying: Why is this that we went out from Egypt?’”
	}
}
{
}

\Verse{21}
{
	\hE{וַיֹּאמֶר מֹשֶׁה שֵׁשׁ מֵאוֹת אֶלֶף רַגְלִי הָעָם אֲשֶׁר אָנֹכִי בְּקִרְבּוֹ וְאַתָּה אָמַרְתָּ בָּשָׂר אֶתֵּן לָהֶם וְאָכְלוּ חֹדֶשׁ יָמִים׃}
}
{
	\eE{
		And Moses said, “The people among whom I am are six
		hundred thousand on foot, and you’ve said, ‘I’ll give them
		meat, and they’ll eat for a month of days!’
	}
}
{
}

\Verse{22}
{
	\hE{הֲצֹאן וּבָקָר יִשָּׁחֵט לָהֶם וּמָצָא לָהֶם אִם אֶת כּל דְּגֵי הַיָּם יֵאָסֵף לָהֶם וּמָצָא לָהֶם׃}
}
{
	\eE{
		Will flock and herd be slaughtered for them, and would it
		provide for them? Or will all the sea’s fish be gathered
		for them, and would it provide for them?”
	}
}
{
}

\Verse{23}
{
	\hE{וַיֹּאמֶר יְהֹוָה אֶל מֹשֶׁה הֲיַד יְהֹוָה תִּקְצָר עַתָּה תִרְאֶה הֲיִקְרְךָ דְבָרִי אִם לֹא׃}
}
{
	\eE{
		And YHWH said to Moses, “Is YHWH’s hand getting short?!
		Now you’ll see if my word happens to you or not.”
	}
}
{
}

\Verse{24}
{
	\hE{וַיֵּצֵא מֹשֶׁה וַיְדַבֵּר אֶל הָעָם אֵת דִּבְרֵי יְהֹוָה וַיֶּאֱסֹף שִׁבְעִים אִישׁ מִזִּקְנֵי הָעָם וַיַּעֲמֵד אֹתָם סְבִיבֹת הָאֹהֶל׃}
}
{
	\eE{
		And Moses went out and spoke YHWH’s words to the people.
		And he gathered seventy men from the people’s elders, and
		he stood them around the tent.
	}
}
{
}

\Verse{25}
{
	\hE{וַיֵּרֶד יְהֹוָה בֶּעָנָן וַיְדַבֵּר אֵלָיו וַיָּאצֶל מִן הָרוּחַ אֲשֶׁר עָלָיו וַיִּתֵּן עַל שִׁבְעִים אִישׁ הַזְּקֵנִים וַיְהִי כְּנוֹחַ עֲלֵיהֶם הָרוּחַ וַיִּתְנַבְּאוּ וְלֹא יָסָפוּ׃}
}
{
	\eE{
		And YHWH came down in a cloud and spoke to him and took
		some of the spirit that was on him and put it on the
		seventy men of the elders. And it was when the spirit
		rested on them: and they prophesied. Then they did not do
		it anymore.
	}
}
{
%	\footnote{
%		they prophesied. What exactly do they do? Prophecy in the Bible does not
%		primarily involve prediction. It is not that they are going about
%		telling the future. What is happening that makes it obvious to Moses,
%		Joshua, and apparently everyone that they are doing something that is
%		associated with prophets? Some suggest that they are in some sort of
%		trance, but trance behavior is not generally part of what is pictured in
%		the fifteen books of the prophets in the Tanak or in the cases of
%		prophecy in the Tanak’s narrative books. What is typical of prophecy, as
%		opposed to historical narrative, in the Bible is that prophecy is in
%		poetry (or in combinations of poetry and prose). Biblical prophecy has
%		the characteristics of oral formulaic poetry. That is, the prophet
%		composes it on the spot, using lines that he or she has already composed
%		and memorized, and mixing these “formulas” with new lines that occur to
%		him or her spontaneously. When such poems occurred to the ancient poets,
%		the poets must have truly felt inspired. And the people who watched and
%		listened to them must have perceived them to be inspired as well. When a
%		man spontaneously says: They’ll beat their swords into plowshares and
%		their spears into pruning hooks. A nation won’t lift a sword against a
%		nation, and they won’t learn war anymore. we can easily imagine him and
%		his audience feeling that it is God moving through him. Further evidence
%		that they are doing something musical or poetic comes from the other
%		most famous biblical instance of such group prophecy: King Saul joins in
%		a group of prophets, and he prophesies with them, which astounds
%		everyone. And in that case the text notes that these prophets are
%		accompanied by musical instruments (1 Sam 10:5). See also the story of
%		Aaron’s and Miriam’s challenge to Moses in the coming chapter. They have
%		experienced prophecy, and they say, “Has YHWH only just spoken through
%		Moses? Hasn’t He also spoken through us?” God responds with a
%		declaration about the nature of prophecy, saying that Moses’ experience
%		is different from all others, and the important thing to note here is
%		that the words of God are in poetry (see Num 12:6–8). Some think that
%		biblical prophecy involved some sort of ecstatic state. That is
%		possible, but it may be more of an inspired state, both in an Isaiah and
%		in a case like Saul’s or the seventy elders here, in which one who does
%		not normally speak in poetry is now moved and enabled to do so.
%	}
}

\Verse{26}
{
	\hE{וַיִּשָּׁאֲרוּ שְׁנֵי אֲנָשִׁים בַּמַּחֲנֶה שֵׁם הָאֶחָד אֶלְדָּד וְשֵׁם הַשֵּׁנִי מֵידָד וַתָּנַח עֲלֵהֶם הָרוּחַ וְהֵמָּה בַּכְּתֻבִים וְלֹא יָצְאוּ הָאֹהֱלָה וַיִּתְנַבְּאוּ בַּמַּחֲנֶה׃}
}
{
	\eE{
		And two men had been left in the camp. One’s name was
		Eldad, and the second’s name was Medad. And the spirit
		rested on them. And they were among the ones who were
		written, but they did not go out to the tent, and they
		prophesied in the camp.
	}
}
{
%	\footnote{
%		they were among the ones who were written. Meaning: they were among the
%		seventy who were written about in the preceding verse.
%	}
}

\Verse{27}
{
	\hE{וַיָּרץ הַנַּעַר וַיַּגֵּד לְמֹשֶׁה וַיֹּאמַר אֶלְדָּד וּמֵידָד מִתְנַבְּאִים בַּמַּחֲנֶה׃}
}
{
	\eE{
		And a boy ran and told Moses, and he said, “Eldad and
		Medad are prophesying in the camp!”
	}
}
{
}

\Verse{28}
{
	\hE{וַיַּעַן יְהוֹשֻׁעַ בִּן נוּן מְשָׁרֵת מֹשֶׁה מִבְּחֻרָיו וַיֹּאמַר אֲדֹנִי מֹשֶׁה כְּלָאֵם׃}
}
{
	\eE{
		And Joshua, son of Nun, Moses’ attendant from his chosen
		ones, answered, and he said, “My lord Moses, restrain
		them!”
	}
}
{
%	\footnote{
%		his chosen ones. The seventy whom Moses had chosen earlier on Jethro’s
%		advice (Exod 18:25). The Aramaic translation (Targum Onkelos) and others
%		read different vowels here and understand this to say: “from his youth.”
%		The Greek (Septuagint) takes it as “chosen,” as does Ibn Ezra. The
%		explicit reference to Moses’ “choosing” a group of seventy in the Jethro
%		episode seems to me to make this the more probable reading here. It also
%		makes Joshua’s concern now more intriguing, as he questions the behavior
%		of members of the second group of seventy. They go beyond what the
%		seventy whom Moses himself had chosen had done, and he wants them to
%		stop. But Moses says, Don’t be jealous for me.
%	}
}

\Verse{29}
{
	\hE{וַיֹּאמֶר לוֹ מֹשֶׁה הַמְקַנֵּא אַתָּה לִי וּמִי יִתֵּן כּל עַם יְהֹוָה נְבִיאִים כִּי יִתֵּן יְהֹוָה אֶת רוּחוֹ עֲלֵיהֶם׃}
}
{
	\eE{
		And Moses said to him, “Are you jealous for me? And who
		would make it so that all of YHWH’s people were prophets,
		that YHWH would put His spirit on them!”
	}
}
{
%	\footnote{
%		who would make it so. Literally, “who would give” or “who would put it.”
%		Meaning: “if only it were so” or “I wish it were so.” This is the only
%		time that Moses uses this expression. The people use it once (Exod
%		16:3); it occurs once as part of a curse (Deut 28:67); and, most
%		extraordinarily, Moses reports that God used it after the revelation at
%		Sinai (Deut 5:29).
%	}
}

\Verse{30}
{
	\hE{וַיֵּאָסֵף מֹשֶׁה אֶל הַמַּחֲנֶה הוּא וְזִקְנֵי יִשְׂרָאֵל׃}
}
{
	\eE{
		And Moses was gathered back to the camp, he and Israel’s
		elders.
	}
}
{
}

\Verse{31}
{
	\hE{וְרוּחַ נָסַע מֵאֵת יְהֹוָה וַיָּגז שַׂלְוִים מִן הַיָּם וַיִּטֹּשׁ עַל הַמַּחֲנֶה כְּדֶרֶךְ יוֹם כֹּה וּכְדֶרֶךְ יוֹם כֹּה סְבִיבוֹת הַמַּחֲנֶה וּכְאַמָּתַיִם עַל פְּנֵי הָאָרֶץ׃}
}
{
	\eE{
		And a wind traveled from YHWH and transported quails from
		the sea and left them on the camp, about a day’s journey
		in one direction and about a day’s journey in the other
		around the camp, and about two cubits on the face of the
		earth.
	}
}
{
%	\footnote{
%		a wind from YHWH. The wording plays exquisitely on the double meaning of
%		the Hebrew word rûa, which means both spirit and wind. (See the
%		comment on Gen 1:2.) Moses has just said, “If only YHWH would put His
%		spirit” on all the people; and now YHWH brings His wind, carrying food
%		for the people. It appears that, in the divine wisdom, the provision of
%		sustenance is of more immediate importance than the bestowing of the
%		spirit. It reminds one of the later rabbinic expression  (’ên
%		qema ’ên tôrh; No flour, no Torah!). Footnote 11:31. two cubits
%		on the face of the earth. Piled three feet high.
%	}
}

\Verse{32}
{
	\hE{וַיָּקם הָעָם כּל הַיּוֹם הַהוּא וְכל הַלַּיְלָה וְכֹל יוֹם הַמּחֳרָת וַיַּאַסְפוּ אֶת הַשְּׂלָו הַמַּמְעִיט אָסַף עֲשָׂרָה חֳמָרִים וַיִּשְׁטְחוּ לָהֶם שָׁטוֹחַ סְבִיבוֹת הַמַּחֲנֶה׃}
}
{
	\eE{
		And the people got up all that day and all night and all
		the next day and gathered the quail. The one who had the
		least gathered ten homers. And they spread them out,
		spreading around the camp.
	}
}
{
%	\footnote{
%		Footnote 11:34. Kibroth Hattaavah. Meaning “Graves of Longing.”
%	}
}

\Verse{33}
{
	\hE{הַבָּשָׂר עוֹדֶנּוּ בֵּין שִׁנֵּיהֶם טֶרֶם יִכָּרֵת וְאַף יְהֹוָה חָרָה בָעָם וַיַּךְ יְהֹוָה בָּעָם מַכָּה רַבָּה מְאֹד׃}
}
{
	\eE{
		And the meat was still between their teeth, not yet
		chewed, and YHWH’s anger flared at the people, and YHWH
		struck at the people, a very great strike.
	}
}
{
}

\Verse{34}
{
	\hE{וַיִּקְרָא אֶת שֵׁם הַמָּקוֹם הַהוּא קִבְרוֹת הַתַּאֲוָה כִּי שָׁם קָבְרוּ אֶת הָעָם הַמִּתְאַוִּים׃}
}
{
	\eE{
		And he called that place’s name Kibroth Hattaavah, because
		they buried the people who were longing there.
	}
}
{
}

\Verse{35}
{
	\hE{מִקִּבְרוֹת הַתַּאֲוָה נָסְעוּ הָעָם חֲצֵרוֹת וַיִּהְיוּ בַּחֲצֵרוֹת׃}
}
{
	\eE{
		From Kibroth Hattaavah the people traveled to Hazeroth,
		and they were in Hazeroth.
	}
}
{
}
